% !PRINTED 296
\weberpage{296}

Ist dagegen $\beta \equiv 0$ und $\alpha-\delta$ nicht durch $p$ teilbar, so
kann man in (9) $\beta'$ beliebig wählen und dann $\alpha'$ aus einer Kongruenz
 ersten Grades bestimmen.

Ist $\beta$ nicht durch $p$ teilbar, so setze man (9) in die Form
\weq{10}{\left[\alpha'\beta+\beta'\left(\frac{\delta-\alpha}{2}\right)\right]^2 \equiv \beta+\beta'^2\left[\left(\frac{\alpha+\delta}{2}\right)^2-1\right] \quad (\mathrm{mod}\ p),}
und wenn nun
\[
\text{II.} \qquad \left(\frac{\alpha+\delta}{2}\right)^2-1 \ \text{quadratischer Rest von}\ p \ \text{ist, so bestimme}
\]
man $\beta'$ aus der Kongruenz:
\[
\beta'^2\left[\left(\frac{\alpha+\delta}{2}\right)^2-1\right] \equiv \left(\frac{\beta-1}{2}\right)^2,
\]
wodurch (10) übergeht in
\[
\alpha'\beta+\beta'\left(\frac{\delta-\alpha}{2}\right) \equiv \pm \frac{\beta+1}{2} \quad (\mathrm{mod}\ p),
\]
und woraus $\alpha'$ bestimmt werden kann, welches Zeichen auch
gewählt wird.

Ist
\[
\text{III.} \qquad \left(\frac{\alpha+\delta}{2}\right)^2-1 \ \text{quadratischer Nichtrest von}\ p \ \text{und zu-}
\]
gleich $\beta$ quadratischer Nichtrest, so läßt sich immer ein Nichtrest
 $\nu$ so bestimmen, daß $\beta+\nu$ quadratischer Rest wird; denn
läßt man $\nu$ in $\beta+\nu$ die Reihe der Nichtreste durchlaufen, so
können nicht \gesperrt{lauter} Nichtreste entstehen, weil unter diesen
auch $\beta$ sein müßte.

Dann kann man $\beta'$ so bestimmen, daß
\[
\beta'^2\left[\left(\frac{\alpha+\delta}{2}\right)^2-1\right] \equiv \nu \quad (\mathrm{mod}\ p),
\]
und dann $\alpha'$ aus
\[
\alpha'\beta+\beta'\left(\frac{\delta-\alpha}{2}\right) \equiv \pm\sqrt{\beta+\nu} \quad (\mathrm{mod}\ p).
\]
Hieraus folgt:

\gesperrt{In einem Normalteiler $\mathfrak{R}$ der Gruppe $\mathfrak{L}_0$ muß gewiß eine
Substitution $U$ von der Form}
\weq{11}{\begin{pmatrix}0, & 1\\ -1, & \xi\end{pmatrix}}
\gesperrt{vorkommen.}

\gesperrt{Es sei zunächst $\xi$ von Null verschieden $(\mathrm{mod}\ p)$.}
